\textbf{Analysis runs in code and visualization in the browser, for researchers and coding agents alike.}
\textbf{a}~Two interfaces over one chromdata.zarr store: the \texttt{uchrom} Python API for analysis, used from notebooks and from code that agents write following skills, and the web browser for visualization, driven by its GUI, by WebMCP tools in the page and from Python through \texttt{uchrom.browser.connect()}, all calling one set of browser actions. Contact maps and multi-modal matrices are linked, not copied (dashed).
\textbf{b}~Linked exploration of the Takei et al.\ 2025 cerebellum data (field of view~0, 47 cells)~\citep{takei2025}, redrawn from the data behind the views: a lasso over the 8 Bergmann glia in the RNA UMAP (1) becomes a group that the 3-D view (2), the Genome view (3; immunofluorescence at the DNA loci and the group's median 3-D distance map, chr15:3--33\,Mb) and the Stats view (4) follow.
\textbf{c}~Top, time from opening the page to the first frame of the 3-D overview (solid) and for the server to open the store (dashed), for the field of view (filled) and copies of its cells tiled 3, 10 and 30 times (open); at most 5\,M points are drawn. Bottom, response time per interaction over HTTP (bars, median of 20 random cells or groups; whiskers, 95th percentile) for the field of view (blue), its tiled copies (circles, 3.8\,M; diamonds, 11.5\,M spots) and 5,313 scHiCAR per-cell contact maps (teal)~\citep{schicar}; dashed line, 100\,ms.
\textbf{d}~Handoff: a researcher selects 8 cells and asks a Claude Code agent (Opus 5.5) what is different about them; following two skills, the agent reads the selection, compares the groups in \texttt{uchrom}, reports the test and caveats and opens the supporting views in the researcher's browser (right, redrawn from the data). The session is shown trimmed (\ldots).
\textbf{e}~Top, the parts of a skill. Bottom, share of grading checks passed on 40 questions with the documentation only (grey) or with the documentation and five skills (orange), by category, overall and on the 20 held-out questions not used while writing the skills. Bars, mean of three models, each question averaged over three runs; whiskers, s.e.m.\ over questions; markers, each model; legend, runs that loaded a skill.
Protocols: Methods.
Logos are marks of their owners, used for identification only.
